\DocumentMetadata{
  lang          = en-US,
  pdfstandard   = {ua-2},
  tagging       = on,
  tagging-setup = {math/setup=mathml-SE} 
}
\documentclass{article}
\def\showanswers{0}
\def\whichlecture{}
\def\lecturetopic{}

% \documentclass{article}
\usepackage{caption}
\usepackage[pdftex]{graphicx}
\usepackage{listings}
\usepackage{tikz}
\usepackage{tikz-qtree}
\usetikzlibrary{arrows,automata}
\usetikzlibrary{shapes.geometric}
\usetikzlibrary{calc,shapes.multipart,chains,arrows}
\usepackage{amsmath}
\usepackage{amssymb}
\usepackage{amsthm}
% \usepackage{enumerate}
% \usepackage{enumitem}
\usepackage{multicol}
% \usepackage{algorithmic}
\usepackage{algorithm,caption}
\usepackage[noend]{algpseudocode}
\usepackage{fancyhdr}
\usepackage{hyperref}
\usepackage{cleveref}
\usepackage[shortlabels]{enumitem}
\usepackage{version}
\usepackage{listings, lstautogobble}
\usepackage{tcolorbox}
\usepackage{wrapfig}

\usepackage{amsmath}
% \usepackage[scaled]{helvet}
% \renewcommand\familydefault{\sfdefault} 
% \usepackage[T1]{fontenc}


\lstset {
    language=C++,
    backgroundcolor=\color{white!5}, % set backgroundcolor
    basicstyle=\footnotesize,% basic font setting
    autogobble = true,
}

\tikzstyle{vertex}=[circle,draw=black,minimum size=20pt,inner sep=0pt]
\tikzstyle{edge} = [draw,thick]
\tikzstyle{weight} = [font=\small]



\theoremstyle{definition}
\newtheorem{exercise}{Exercise}%[section]
\newtheorem{question}{Question}%[section]


%% For use with balanced BSTs
\tikzset{
itria/.style={
  draw,dashed,shape border uses incircle,
  isosceles triangle,shape border rotate=90,yshift=-1.45cm}}

\newenvironment {solution}
{
\begin{tcolorbox}
}
{
\end{tcolorbox}
}



\newcommand{\qacc}{q_{\textnormal{accept}}}
\newcommand{\qrej}{q_{\textnormal{reject}}}



\title{Unit 2 Group Problem Session}
\author{
    Daniel Frishberg \\
    Chris Amey, Alex Hanus, Zachary Harder, Nolan Knievel, and Ronan Biggs    
}
\date{Cal Poly CSC 3445, Fall 2026}

\begin{document}
\lstset{upquote=true}
\setlength\parindent{0em}

\maketitle

\section{Instructions}
Solve these problems together, in a group of four.

Each team member will write up the solution to one of the questions.

Submit your solutions to Gradescope (more instructions to follow in Canvas). You will need to provide:

\begin{itemize}
    \item The names of all team members
    \item Your solutions to the problems, indicating who wrote up each solution.
    \item \textbf{You may consult generative AI tools and other sources, but do not use them to write your solutions.}
    \item Cite any external sources or tools used and you used them, or indicate that you did not use any. \textbf{This part is required.}
\end{itemize}

\pagebreak

\section{Questions}
\begin{enumerate}
\item Let $\Sigma = \{a, b, c\}$. Draw a PDA that recognizes the following language:
\[
    \{ a^kb^{2l}c^m \mid l \geq 0, m > k \geq 0\} 
\]


\item Write a CFG that generates the language
\[
    \{a^kb^{2l}c^m \mid l \geq k \geq 0, m \geq 0\}
\]
\textbf{Note this is NOT the language in the previous question.}


The CFG does \textbf{not} have to be in Chomsky normal form (CNF).

\item Consider the following CFG.
\begin{align*}
    S &\rightarrow ZT \mid YR \mid ZZ \mid YY \mid 0 \mid 1\\
    X &\rightarrow ZT \mid YR \mid ZZ \mid YY \mid 0 \mid 1\\   
    R &\rightarrow XY\\
    T &\rightarrow XZ \\
    Z &\rightarrow 0 \\
    Y &\rightarrow 1
\end{align*} 

Consider the string $00100$. Run the CYK algorithm to determine whether $G$ generates this string. If it does, recover the parse tree that results from the algorithm. If not, say why not.

\item Consider the language $L = \{r\#st\#ts \mid r,s,t \in \{0,1\}^* \}$. Prove that $L$ is not context free using the pumping lemma for context-free languages.

\end{enumerate}
\end{document}